\documentclass[../main.tex]{subfiles}

\begin{document}

\section{Chapter 6: Visual Integrity (Ink, Overlap, Color, Layout, and Annotation)}
\label{sec:chapter6-visual-integrity}

\subsection{Techniques / Principles Applied}

\begin{itemize}[leftmargin=*]
\item Proportional ink and true baselines for quantitative marks.
\item Color choices that preserve contrast and basic accessibility.
\item Overlap control through structural separation, spacing, and layout.
\item Data--context balance and annotation grammar that clarify interpretation
without adding chartjunk.
\end{itemize}

\subsection{How Applied in the Dashboard}

\subsubsection{Proportional Ink and True Baselines}

The dashboard maintains proportional ink by anchoring bar charts at true zero
baselines and by keeping bar length or flow width tied to the underlying
weighted values. This is visible in
\figref{fig:cash-only-reasons}{Cash-only reasons} and in
\figref{fig:account-composition-flow}{Account composition flow}.

\subsubsection{Color and Accessibility}

Color is used with restraint.
\figref{fig:socioeconomic-matrix}{Socioeconomic matrix} keeps a single
sequential palette so density is shown through intensity rather than extra hue
variation, while \figref{fig:national-inflows}{National inflow digitization}
uses a limited set of contrasting fills to separate payment channels without
turning the page into a saturated categorical display.

\subsubsection{Managing Overplotting and Graphical Overlap}

To keep dense comparisons readable, the dashboard separates marks instead of
forcing too many categories into a single crowded frame.

\dashboardfigure{0.78}{fig-page3-exclusion-friction-log.png}
{The Exclusion Friction Log: gap analysis benchmarking unbanked constraints vs. national baseline.}
{fig:exclusion-friction-log}

\figref{fig:mobile-ownership-profile}{Mobile ownership profile} splits the
ownership analysis into small multiples, so each subgroup comparison stays
scannable rather than collapsing into overlapping grouped bars.
\figref{fig:national-inflows}{National inflow digitization} and
\figref{fig:exclusion-friction-log}{Exclusion friction log} both use generous
row spacing, which prevents labels, benchmark markers, and adjacent bars from
colliding. On Page 3, \figref{fig:account-composition-flow}{Account composition flow}
also reduces visual clutter by grouping and ordering ribbons into readable
streams instead of a tangled crossover pattern.

\subsubsection{Data--Context Balance and Annotation Grammar}

Data marks are not presented in isolation. Titles, subtitles, legends, and
benchmark cues explain what is being compared and how the reader should
interpret the result without competing with the data itself.
\figref{fig:exclusion-friction-log}{Exclusion friction log} is the clearest
example: the benchmark line, legend, and direct barrier labels let the reader
compare a filtered subgroup against the national rate in the same frame, so
context supports the analysis instead of overwhelming it.

\subsection{Notes / Adjustments}

\figref{fig:national-inflows}{National inflow digitization} uses a 100\%
stacked format, so only the first segment in each row has a common baseline.
That reduces precision for the inner segments, but the compromise is
acceptable because the page is intended to show composition rather than exact
segment-by-segment ranking.

\end{document}
